\newpage
\phantomsection
\label{prf:euclid-i-4}
\LRAProofFor{thm:euclid-i-4}

\begin{remark*}[Return]
\hyperref[thm:euclid-i-4]{Return to Theorem}
\end{remark*}

\begin{theorem*}[Euclid I.4: Side-angle-side congruence]
If two triangles have two sides equal to two sides respectively, and have the
angles contained by the equal straight lines equal, then the bases, the
triangles, and the remaining angles are equal respectively.
\end{theorem*}

\begin{center}
\begin{tikzpicture}[scale=0.95, line cap=round, line join=round]
  \definecolor{euclidblue}{RGB}{30,110,190}
  \tikzset{
    onetick/.style={postaction={decorate},decoration={markings,mark=at position 0.5 with {\draw (0,-2pt)--(0,2pt);}}},
    twotick/.style={postaction={decorate},decoration={markings,mark=at position 0.45 with {\draw (0,-2pt)--(0,2pt);},mark=at position 0.55 with {\draw (0,-2pt)--(0,2pt);}}}
  }
  \coordinate (A) at (0,0);
  \coordinate (B) at (2.6,0);
  \coordinate (C) at (1.0,1.8);
  \coordinate (D) at (4.3,0);
  \coordinate (E) at (6.9,0);
  \coordinate (F) at (5.3,1.8);
  \draw[very thick, onetick] (A)--(B);
  \draw[very thick, twotick] (A)--(C);
  \draw[very thick, euclidblue] (B)--(C);
  \draw[very thick, onetick] (D)--(E);
  \draw[very thick, twotick] (D)--(F);
  \draw[very thick, euclidblue] (E)--(F);
  \draw (0.38,0) arc[start angle=0,end angle=61,radius=0.38];
  \draw ($(D)+(0.38,0)$) arc[start angle=0,end angle=61,radius=0.38];
  \foreach \p/\pos in {A/below left,B/below,C/above,D/below left,E/below,F/above}{
    \fill (\p) circle (0.026);
    \node[\pos] at (\p) {$\p$};
  }
  \node[font=\footnotesize] at (1.35,-0.45) {$AB=DE$};
  \node[font=\footnotesize] at (5.62,-0.45) {$AC=DF$};
  \node[euclidblue, font=\footnotesize] at (3.45,1.15) {included angles equal};
\end{tikzpicture}
\end{center}


\begin{proof}[Professional Standard Proof]
\LRAProofBodyStart
Let triangles $ABC$ and $DEF$ satisfy $AB=DE$, $AC=DF$, and
$\angle BAC=\angle EDF$. Place the point $A$ on $D$, the straight line $AB$ on
$DE$, and the straight line $AC$ on $DF$. Since the included angles are equal,
the two sides fall along the corresponding sides. Since $AB=DE$, the point $B$
coincides with $E$; since $AC=DF$, the point $C$ coincides with $F$.

Therefore the base $BC$ coincides with the base $EF$. By Common Notion IV,
things which coincide are equal. Hence $BC=EF$, the whole triangle $ABC$ is
equal to the whole triangle $DEF$, and the remaining corresponding angles are
equal.
\end{proof}

\begin{proof}[Detailed Learning Proof]
\LRAProofBodyStart
The diagram is meant to be read by superposition. Equal included angles allow
one triangle to be laid on the other without rotating one side away from its
mate. The two equal side lengths then force the two non-vertex endpoints to
land exactly on the corresponding endpoints. Once the three vertices coincide,
the third side and the two remaining angles coincide as well.
\end{proof}

\begin{remark*}[Proof structure]
Superpose one triangle on the other, use the two equal sides to identify the
endpoints, and use coincidence to conclude equality of the base and remaining
angles.
\end{remark*}

\begin{dependencies}
\begin{itemize}
  \item \hyperref[def:euclid-triangle-side-classes]{Triangle side classes}
  \item \hyperref[def:euclid-rectilinear-angle]{Rectilinear angle}
  \item \hyperref[ax:euclid-common-notion-coinciding-things-equal]{Euclid Common Notion IV}
\end{itemize}
\end{dependencies}

\clearpage
